\documentclass[twocolumn]{aastex631}

\usepackage{booktabs}
\usepackage{amsmath}
\usepackage{graphicx}

\shorttitle{Calibration of Production Transient Classifiers}
\shortauthors{Sati}

\begin{document}

\title{Calibration of Production Transient Classifiers: A Benchmark Study}

\author{Pallavi Sati}
\affiliation{University of Bristol, School of Physics, Bristol BS8 1TL, UK}

\begin{abstract}
The Vera C.\ Rubin Observatory's Legacy Survey of Space and Time will 
generate $\sim$$10^7$ transient alerts per night against a global 
spectroscopic follow-up capacity of $\sim$$10^5$ observations per year. 
Automated alert brokers address this mismatch by outputting class 
probability vectors, which astronomers use to prioritise follow-up 
observations. These probabilities are implicitly treated as reliable 
confidence estimates, yet no published work has evaluated the 
calibration of production transient classifiers against independent 
spectroscopic ground truth.

We present the first systematic calibration audit of production 
astronomical transient classifiers, evaluating ALeRCE's light curve 
classifier and two Fink classifiers — a Random Forest and SuperNNova 
binary discriminator — on 1,606 spectroscopically confirmed objects 
from the ZTF Bright Transient Survey. We compute Expected Calibration 
Error (ECE) with bootstrap confidence intervals and visualise 
miscalibration via reliability diagrams.

ALeRCE's classifier exhibits systematic underconfidence 
($\mathrm{ECE} = 0.259$; 95\% CI: $[0.237, 0.280]$), with predicted 
probabilities consistently below empirical accuracies across all 
confidence bins. Miscalibration follows an inverse relationship with 
classification accuracy: SNIa and SNIbc, classified with accuracies 
above 0.83, exhibit the worst calibration (ECE $> 0.36$), while SNII, 
classified at 0.495 accuracy, is the best-calibrated class 
(ECE $= 0.103$). This pattern is mechanistically explained by Random 
Forest probability suppression under class imbalance. TDE achieves 
classification accuracy of 0.000 due to a taxonomy gap in ALeRCE's 
output space — a finding with direct operational consequences for 
LSST-era science. Temperature scaling with a single scalar parameter 
($T = 0.359$) reduces ALeRCE's aggregate ECE by 88\% to $0.031$, 
requiring no model retraining.

Fink's classifiers exhibit structural miscalibration 
($\mathrm{ECE} = 0.464$ for RF, $0.304$ for SuperNNova) that 
post-hoc calibration methods cannot adequately address; temperature 
scaling optimisation hits the upper bound for both, indicating 
non-monotonic score distributions incompatible with standard 
calibration techniques.

These results demonstrate that miscalibration is systematic, 
broker-dependent, and consequential for confidence-based 
follow-up prioritisation. Post-hoc calibration can recover 
reliable probabilities for ALeRCE but not for Fink's current 
classifiers. We argue that calibration metrics should be 
reported alongside accuracy for all production astronomical 
classifiers, and that the calibrated ALeRCE outputs provide 
a suitable foundation for deferral-aware classification 
frameworks targeting LSST-era spectroscopic resource allocation.
\end{abstract}

\keywords{astronomical databases --- surveys --- supernovae: general --- methods: statistical}

\section{Introduction}
\label{sec:intro}
The Vera C.\ Rubin Observatory's Legacy Survey of Space and Time 
\citep[LSST;][]{Ivezic2019} will generate approximately $10^7$ 
transient alerts per night beginning in 2025. Global spectroscopic 
follow-up capacity — the essential confirmation mechanism for 
transient classification — is fixed at roughly $10^5$ observations 
per year \citep{Hambly2022}. This mismatch of four orders of magnitude 
transforms transient astronomy from a data-limited to a 
decision-limited discipline: the central challenge is no longer 
detecting transients but determining which ones warrant the scarce 
resource of spectroscopic confirmation.

Automated alert brokers — ALeRCE \citep{Forster2021}, Fink 
\citep{Moller2021}, ANTARES \citep{Narayan2018}, and Lasair 
\citep{Smith2019} — address this challenge by classifying 
transients from photometric light curves and outputting class 
probability vectors. These probabilities are used by astronomers 
and automated pipelines to prioritise follow-up observations. 
Implicit in this usage is an assumption that the probabilities 
are meaningful confidence estimates — that a classifier outputting 
$P(\mathrm{SNIa}) = 0.9$ is correct approximately 90\% of the time.

This assumption has never been tested. No published work evaluates 
the calibration of production transient classifiers using 
independent spectroscopic ground truth. Reliability diagrams and 
Expected Calibration Error \citep[ECE;][]{Naeini2015} — standard 
diagnostic tools in machine learning — are absent from the 
astronomical broker literature. If production classifiers are 
systematically miscalibrated, the confidence-based prioritisation 
decisions made by astronomers and automated systems are built on 
unreliable foundations.

In this paper we conduct the first systematic calibration audit 
of production transient classifiers, evaluating ALeRCE's light 
curve classifier and two Fink classifiers on 1,606 
spectroscopically confirmed objects from the ZTF Bright Transient 
Survey \citep{Fremling2020, Perley2020}. Our principal findings are:

\begin{enumerate}
\item ALeRCE's classifier exhibits systematic underconfidence 
($\mathrm{ECE} = 0.259$), not the overconfidence typical of 
deep learning models. Predicted probabilities are consistently 
lower than empirical accuracies across all confidence bins.

\item The miscalibration follows an inverse relationship with 
accuracy: the best-classified types (SNIa, SNIbc; accuracy $>0.83$) 
are the most severely underconfident (ECE $>0.36$), while the 
worst-classified type (SNII; accuracy $0.495$) is the 
best-calibrated (ECE $0.103$).

\item Temperature scaling with a single scalar parameter 
($T = 0.359$) reduces ALeRCE's ECE by 88\%, to $0.031$ — 
below the threshold for practically significant miscalibration, 
with no model retraining required.

\item TDE, a high-priority LSST science target, achieves 
classification accuracy of $0.000$ on ALeRCE because TDE 
is absent from ALeRCE's output taxonomy.

\item Fink's classifiers exhibit structural miscalibration 
($\mathrm{ECE} = 0.464$ for RF, $0.304$ for SuperNNova) 
that post-hoc calibration methods cannot adequately address.
\end{enumerate}

These findings have immediate implications for spectroscopic 
follow-up prioritisation and motivate the deferral-aware 
classification framework we develop in subsequent work. 
Section~\ref{sec:related} reviews related work on broker 
classification and calibration methods. 
Section~\ref{sec:data} describes our dataset. 
Section~\ref{sec:methods} details our calibration methodology. 
Section~\ref{sec:results} presents results. 
Section~\ref{sec:discussion} discusses implications and limitations.


\section{Related Work}
\label{sec:related}
% Written after results

\section{Data}
\label{sec:data}
\subsection{The ZTF Bright Transient Survey}
\label{sec:data:bts}

Our dataset is drawn from the Zwicky Transient Facility Bright Transient 
Survey \citep[BTS;][]{Fremling2020, Perley2020}, a magnitude-limited 
($m < 18.5$) spectroscopic completeness survey operating since 2018. 
The BTS provides spectroscopic classifications for ZTF transients, 
making it the largest homogeneous sample of spectroscopically confirmed 
extragalactic transients currently available. We use spectroscopic 
classifications as ground truth labels throughout this work.

\subsection{Sample Construction}
\label{sec:data:sample}

We queried the ALeRCE broker API \citep{Forster2021} for classification 
probabilities corresponding to all BTS objects with confirmed 
spectroscopic types on 22 February 2026. ALeRCE's light curve classifier 
\citep{Sanchez-Saez2021} outputs probabilities across 15 classes derived 
from features of the ZTF alert photometric time series. For each object 
we retrieved the most recent classification available at query time, 
corresponding to classifier versions \texttt{24.4.1}--\texttt{27.5.7a32.dev1} 
active in the production pipeline. We provide the full list of ZTF object 
identifiers and retrieved probability vectors as a machine-readable 
supplement to this paper.

We map the 11 BTS spectroscopic types to the 5 ALeRCE transient classes 
most relevant to extragalactic time-domain science: Type Ia supernovae 
(SNIa), Type II supernovae (SNII), stripped-envelope supernovae (SNIbc), 
superluminous supernovae (SLSN), and tidal disruption events (TDE). 
The full mapping is given in Table~\ref{tab:mapping}.

Of 1,987 BTS objects queried, 1,606 returned valid ALeRCE classifications 
(80.1\% success rate). Failed queries reflect objects observed before 
ALeRCE's alert ingestion began or with insufficient photometric history 
for classification. The final sample comprises 757 SNIa, 560 SNII, 
204 SNIbc, 55 SLSN, and 30 TDE.

To avoid class imbalance biasing our rarest-class estimates, we collected 
all available BTS objects in the SLSN and TDE categories rather than 
applying a fixed per-class cap. This yields a sample in which common 
classes (SNIa, SNII) are well-represented while preserving the full 
available statistics for scientifically rare but high-priority classes.

\begin{deluxetable}{ll}
\tablecaption{BTS spectroscopic type to ALeRCE class mapping.
\label{tab:mapping}}
\tablehead{\colhead{BTS Types} & \colhead{ALeRCE Class}}
\startdata
SN~Ia, SN~Ia-91T, SN~Ia-91bg, SN~Ia-pec & SNIa \\
SN~II, SN~IIn, SN~IIb, SN~IIP           & SNII \\
SN~Ib, SN~Ic, SN~Ib/c, SN~Ic-BL         & SNIbc \\
SLSN-I                                   & SLSN \\
TDE                                      & TDE \\
\enddata
\end{deluxetable}

\subsection{Fink Broker Data}
\label{sec:data:fink}

We additionally queried the Fink broker \citep{Moller2021} for the same 
set of ZTF objects via the Fink Portal API on 24 February 2026, using 
HTTP POST requests to \texttt{api.fink-portal.org/api/v1/objects}. 
Fink provides two binary classifiers relevant to our analysis. 
\texttt{snn\_snia\_vs\_nonia} is based on the SuperNNova deep learning 
classifier \citep{Moller2019} and outputs a general $P(\mathrm{SNIa})$ 
score. \texttt{rf\_snia\_vs\_nonia} is a Random Forest classifier 
documented as targeting \textit{rising} SN Ia candidates; its scores 
are designed for early-phase, pre-maximum alerts rather than 
general-purpose SN Ia posteriors. We include both scores in our 
analysis but explicitly flag this distinction when interpreting the 
RF results in Sections~\ref{sec:results:fink} and~\ref{sec:discussion:fink}.

For each object we retrieved the most recent alert-level score available 
in the Fink database. We distinguish three outcomes: object found with 
scores (collected), object absent from Fink (not in ingestion window, 
treated as missing data), and API error (excluded). Absent objects are 
not coerced to zero — they are excluded from the Fink analysis entirely.

Of 1,606 objects queried, 1,368 (85.2\%) were present in the Fink 
database, with 237 objects absent due to falling outside Fink's alert 
ingestion window. Both classifiers returned scores for all 1,368 matched 
objects, yielding a comparative subsample of 644 SNIa, 480 SNII, 
166 SNIbc, 50 SLSN, and 28 TDE.

\subsection{Ground Truth and Selection Effects}
\label{sec:data:groundtruth}

We emphasise that our ground truth labels derive from independent 
spectroscopic observations, not from the broker classifications being 
evaluated. This avoids the circular validation present in studies that 
use broker-assigned labels to assess broker performance.

The BTS magnitude limit ($m < 18.5$) introduces a selection bias toward 
brighter, lower-redshift transients. Our sample therefore represents 
the regime in which spectroscopic follow-up is most feasible, which is 
precisely the regime most relevant to the deferral decisions this work 
motivates. We do not claim our calibration results generalise to 
fainter or higher-redshift populations. As the BTS selects transients 
for spectroscopic follow-up partly on the basis of brightness and 
observability, our sample likely over-represents well-behaved, 
unambiguous examples of each class. Calibration performance on the 
fainter, more ambiguous objects that will dominate LSST alert streams 
may differ; our results should be interpreted as an upper bound on 
broker performance in the bright regime.

\section{Methods}
\label{sec:methods}
\subsection{Expected Calibration Error}
\label{sec:methods:ece}

A classifier is said to be \textit{calibrated} if its predicted 
confidence matches its empirical accuracy. Formally, for a classifier 
outputting confidence $\hat{p}$ and true label $y$:

\begin{equation}
P(y = \hat{y} \mid \hat{p} = p) = p \quad \forall\, p \in [0, 1]
\label{eq:calibration}
\end{equation}

where $\hat{y}$ is the predicted class. We measure calibration using 
Expected Calibration Error \citep[ECE;][]{Naeini2015}, which 
approximates the expectation in Equation~\ref{eq:calibration} by 
partitioning predictions into $M$ equally spaced confidence bins:

\begin{equation}
\mathrm{ECE} = \sum_{m=1}^{M} \frac{|B_m|}{n} 
\left| \mathrm{acc}(B_m) - \mathrm{conf}(B_m) \right|
\label{eq:ece}
\end{equation}

where $B_m$ is the set of samples whose top-1 confidence falls in bin 
$m$, $n$ is the total number of samples, $\mathrm{acc}(B_m)$ is the 
fraction of correct predictions in bin $m$, and $\mathrm{conf}(B_m)$ 
is the mean confidence in bin $m$. We use $M = 10$ bins throughout. 
A perfectly calibrated classifier has $\mathrm{ECE} = 0$; following
common practice we use 0.05 as a practical reference level for
consequential miscalibration \citep{Guo2017}, while noting this is a
convention, not a formal significance boundary.

For ALeRCE's 15-class classifier, we compute ECE using the top-1
predicted class and its associated probability across all 15 output
classes. For Fink's binary classifiers, which output a scalar
$P(\mathrm{SNIa})$, we apply a binary variant of
Equation~\ref{eq:ece} where $\mathrm{acc}(B_m)$ is the observed
SNIa fraction within bin $m$.

\subsubsection{Per-class calibration: three definitions}
\label{sec:methods:perclass}

``Per-class calibration'' is ambiguous in the multiclass setting, and
the choice of definition materially affects conclusions
(Section~\ref{sec:results:alerce_perclass}). We therefore compute and
report three distinct quantities explicitly:

\begin{enumerate}
\item \textbf{True-class-conditioned top-label ECE}: for class $k$,
Equation~\ref{eq:ece} restricted to objects whose \emph{spectroscopic}
label is $k$. This measures a class-conditional confidence deficit —
how appropriately the classifier weights objects that truly are class
$k$ — which is the operationally relevant quantity when asking whether
genuine members of a class will clear a confidence-based follow-up
threshold. It is \emph{not} calibration in the standard conditional
sense.
\item \textbf{Predicted-class-stratified top-label ECE}:
Equation~\ref{eq:ece} restricted to objects whose \emph{predicted}
top-1 class is $k$. This asks whether confidence matches precision
within each predicted-class stream.
\item \textbf{One-vs-rest classwise ECE}: over \emph{all} objects,
binning the class-$k$ probability $p_k$ against the empirical
frequency of true class $k$ — the standard classwise calibration
condition $P(y = k \mid \hat{p}_k = p) = p$.
\end{enumerate}

\subsubsection{Ranking metrics}
\label{sec:methods:ranking}

ECE alone cannot distinguish a classifier that ranks objects usefully
but reports distorted probabilities (repairable post hoc) from one
with no discriminative signal (not repairable by any monotone
transformation, since monotone recalibration preserves ranking). For
the binary Fink classifiers we therefore also report ROC--AUC and
average precision (PR--AUC), alongside Brier score. A score with
$\mathrm{AUC} \approx 0.5$ can always be made to \emph{appear}
calibrated — a constant prediction at the base rate has
$\mathrm{ECE} \approx 0$ — while carrying no information; reporting
ECE together with ranking metrics guards against this failure mode.

\subsubsection{Robustness to bin construction}
\label{sec:methods:ace}

Equation~\ref{eq:ece}'s equal-width bins are known to be sensitive
to how probability mass is distributed across the confidence range
\citep{Nixon2019, Roelofs2022}: a classifier whose scores concentrate
near one end of $[0,1]$ — Fink's \texttt{rf\_snia\_vs\_nonia} has mean
score $0.011$, putting roughly 90\% of objects in the first of 10
bins — leaves the remaining bins populated by very few points, so the
resulting ECE can be more sensitive to noise or bin placement than a
well-populated case. As a robustness check, we additionally compute
each headline ECE using equal-\textit{mass} (adaptive) bins, i.e.
quantile bins of the confidence/score distribution that guarantee
$\approx n/10$ samples per bin regardless of skew. If equal-width and
equal-mass binning agree, the reported miscalibration is a property
of the classifier and not an artefact of bin construction.

\subsection{Uncertainty Quantification via Bootstrap}
\label{sec:methods:bootstrap}

Point estimates of ECE carry substantial sampling uncertainty, 
particularly for small classes such as SLSN ($n=55$) and TDE ($n=30$). 
We report 95\% confidence intervals on all ECE estimates using the 
nonparametric bootstrap \citep{Efron1979}: we resample the evaluation 
set with replacement $B = 2000$ times, compute ECE on each resample, 
and report the 2.5th and 97.5th percentiles of the resulting 
distribution. All bootstrap procedures use a fixed random seed 
(seed $= 42$) for reproducibility.

\subsection{Reliability Diagrams}
\label{sec:methods:reliability}

We visualise calibration using reliability diagrams \citep{DeGroot1983}, 
which plot observed accuracy against mean predicted confidence for each 
bin. A perfectly calibrated classifier produces points lying on the 
diagonal. Points above the diagonal indicate underconfidence 
(predicted probability lower than empirical accuracy); points below 
indicate overconfidence. We include sample-count histograms alongside 
each diagram to indicate where estimates are statistically reliable.

\subsection{Temperature Scaling}
\label{sec:methods:tempscaling}

Temperature scaling \citep{Guo2017} is a post-hoc calibration method 
that applies a single scalar parameter $T > 0$ to the pre-softmax 
logits of a classifier. For a classifier outputting probability vector 
$\mathbf{p}$, the calibrated probabilities are:

\begin{equation}
\hat{\mathbf{p}}_T = \mathrm{softmax}\!\left(\frac{\log \mathbf{p}}{T}\right)
\label{eq:tempscaling}
\end{equation}

Values of $T < 1$ increase confidence (sharpening the distribution); 
values of $T > 1$ decrease confidence (flattening toward uniform). 
The optimal $T$ is found by minimising ECE on the evaluation set 
using bounded scalar optimisation over $T \in [0.1, 10.0]$.

For Fink's binary classifiers, we apply an equivalent logit-space 
temperature scaling:

\begin{equation}
\hat{p}_T = \sigma\!\left(\frac{\mathrm{logit}(p)}{T}\right)
\label{eq:tempscaling_binary}
\end{equation}

where $\sigma$ is the sigmoid function and $\mathrm{logit}(p) = 
\log(p / (1-p))$.

Fitting $T$ on the same data used to evaluate it yields an optimistic
estimate. Our primary results therefore use held-out evaluation
throughout: 5-fold stratified cross-validation ($T$ fitted on four
folds by NLL minimisation, ECE evaluated on the held-out fold, pooled
across folds), corroborated by an independent 60/40
calibration/test split. We additionally report the same-set fit as an
explicit \emph{oracle upper bound}, labelled as such, for comparison
with the held-out numbers.

For the Fink classifiers we also fit isotonic regression
\citep{Zadrozny2002} under the same 5-fold cross-validation — a
substantially more flexible monotone calibrator than scalar
temperature — so that failure of temperature scaling alone is not
over-interpreted as failure of all post-hoc methods. Isotonic results
must be read jointly with the ranking metrics of
Section~\ref{sec:methods:ranking}: a calibrator applied to an
uninformative score can achieve low ECE by collapsing predictions
toward the base rate.

\subsection{Probability Renormalisation}
\label{sec:methods:renorm}

ALeRCE's classifier outputs probabilities across 15 classes, of which 
we evaluate 4 (excluding TDE, which is absent from ALeRCE's taxonomy). 
To assess whether apparent miscalibration arises from probability 
dilution across unevaluated classes, we compute a renormalised ECE 
by projecting the 15-class probability vector onto the 4 evaluated 
classes and renormalising to sum to unity. The difference between 
15-class and renormalised ECE quantifies the contribution of probability 
dilution to the observed miscalibration.

\section{Results}
\label{sec:results}
\subsection{ALeRCE Aggregate Calibration}
\label{sec:results:alerce_agg}

Figure~\ref{fig:alerce_reliability} shows the reliability diagram for 
ALeRCE's light curve classifier evaluated on 1,576 spectroscopically 
confirmed objects (30 TDE excluded; see Section~\ref{sec:results:tde}). 
The classifier exhibits systematic underconfidence across all confidence 
bins: every point lies above the perfect calibration diagonal, indicating 
that predicted probabilities are consistently lower than empirical 
accuracies.

\begin{figure}
\centering
\includegraphics[width=\columnwidth]{../figures/fig1_reliability_diagram}
\caption{Reliability diagram for ALeRCE's light curve classifier 
evaluated on 1,576 spectroscopically confirmed ZTF transients 
(30 TDE excluded). Every point lies above the perfect calibration 
diagonal, indicating systematic underconfidence across all confidence 
bins. The histogram shows sample counts per bin. 
$\mathrm{ECE} = 0.259$ (95\% CI: $[0.237, 0.280]$).}
\label{fig:alerce_reliability}
\end{figure}

The aggregate ECE is $0.259$ (95\% CI: $[0.237, 0.280]$), well above 
the 0.05 threshold for practically significant miscalibration. The 
overall accuracy is $0.709$ while the mean top-1 confidence is $0.451$, 
a gap of $+0.258$ that is consistent with the ECE estimate. This gap 
is directional and consistent — not a consequence of random binning 
artefacts — indicating a systematic mechanistic suppression of 
classifier confidence.

To assess whether this miscalibration arises from probability dilution 
across ALeRCE's 15-class output space, we renormalise the probability 
vectors onto the 4 evaluated classes and recompute ECE. The renormalised 
ECE is $0.251$ (95\% CI: $[0.230, 0.272]$), a difference of $0.008$ 
from the 15-class estimate. Probability dilution accounts for only 
3\% of the observed miscalibration; the remainder is intrinsic to the 
classifier.

\subsection{ALeRCE Per-Class Calibration}
\label{sec:results:alerce_perclass}

Figure~\ref{fig:alerce_perclass} shows per-class reliability diagrams 
for the four ALeRCE-classified transient types. The results reveal a 
striking inverse relationship between classification accuracy and 
calibration error.

SNIa, the most accurately classified class (accuracy $= 0.831$), 
exhibits the worst calibration ($\mathrm{ECE} = 0.377$, 
95\% CI: $[0.352, 0.401]$). SNIbc shows the same pattern 
(accuracy $= 0.828$, $\mathrm{ECE} = 0.362$, 
95\% CI: $[0.315, 0.405]$). Both classes are classified correctly 
more than 80\% of the time, yet the classifier assigns them 
confidences far below this accuracy level.

SNII presents the inverse case: accuracy is only $0.495$, yet ECE 
is $0.103$ (95\% CI: $[0.066, 0.141]$) — the best-calibrated class. 
The classifier is appropriately uncertain about objects it 
classifies poorly. SLSN shows intermediate behaviour 
($\mathrm{ECE} = 0.245$, 95\% CI: $[0.181, 0.319]$), with wider 
confidence intervals reflecting the smaller sample size ($n = 55$).

This pattern is mechanistically coherent. ALeRCE's light curve 
classifier is based on Random Forest \citep{Sanchez-Saez2021}, 
whose probability outputs are averages of hard leaf-node votes. 
In class-imbalanced datasets, this averaging systematically 
suppresses predicted probabilities for well-represented classes 
\citep{Niculescu-Mizil2005, Zadrozny2002}. SNIa comprises 47\% 
of our sample; its dominance mechanically drives down predicted 
probabilities across the board, with the strongest suppression 
for the class the classifier discriminates most reliably.

\begin{figure}
\centering
\includegraphics[width=\columnwidth]{../figures/fig2_perclass_reliability}
\caption{Per-class reliability diagrams for ALeRCE's light curve 
classifier. SNIa and SNIbc (accuracy $> 0.83$) show severe 
underconfidence (ECE $> 0.36$). SNII (accuracy $0.495$) is the 
best-calibrated class (ECE $= 0.103$). SLSN shows wider confidence 
intervals reflecting the smaller sample size ($n = 55$).}
\label{fig:alerce_perclass}
\end{figure}

\subsection{TDE Classification Failure}
\label{sec:results:tde}

TDE achieves a classification accuracy of $0.000$ on our 30-object 
sample. This is not a calibration problem but a taxonomy gap: TDE 
is absent from ALeRCE's 15-class output taxonomy. The classifier 
cannot predict TDE by design, assigning all probability mass to 
other classes for every TDE in our sample.

This finding has direct operational significance. TDEs are a 
high-priority science case for LSST-era time-domain astronomy, 
representing unique probes of supermassive black hole demographics 
and accretion physics \citep{Gezari2021}. A production broker 
used by the community cannot identify them. We exclude TDE from 
all calibration analyses but flag this gap as requiring attention 
before LSST first light.

\subsection{Temperature Scaling}
\label{sec:results:tempscaling}

We apply temperature scaling (Section~\ref{sec:methods:tempscaling}) 
to ALeRCE's classifier outputs. The optimal temperature is 
$T = 0.359$, substantially below unity, confirming that the 
classifier requires confidence \textit{sharpening} — the opposite 
of the overconfidence correction for which temperature scaling 
is most commonly applied \citep{Guo2017}.

Figure~\ref{fig:tempscaling} shows the reliability diagram before 
and after scaling. Post-calibration ECE is $0.031$ 
(95\% CI: $[0.017, 0.047]$), an 88\% reduction from the 
pre-calibration value of $0.259$. The confidence intervals 
do not overlap, confirming a statistically significant improvement. 
Mean confidence increases from $0.451$ to $0.712$, now closely 
matching the empirical accuracy of $0.709$.

\begin{figure}
\centering
\includegraphics[width=\columnwidth]{../figures/fig4_temperature_scaling}
\caption{ALeRCE reliability diagrams before (left) and after (right) 
temperature scaling with $T = 0.359$. ECE reduces from $0.259$ to 
$0.031$ — an 88\% improvement. The residual dip near confidence 
$0.3$ in the right panel reflects SNII overcorrection 
(see Section~\ref{sec:discussion:tempscaling}).}
\label{fig:tempscaling}
\end{figure}

Table~\ref{tab:tempscaling} shows per-class ECE before and after 
scaling. Temperature scaling substantially improves SNIa 
($0.377 \to 0.114$) and SNIbc ($0.360 \to 0.108$), with moderate 
improvement for SLSN ($0.225 \to 0.140$). However, SNII 
\textit{worsens} from $0.099$ to $0.189$. 

This result is mechanistically interpretable. SNII was the only 
well-calibrated class before scaling: the classifier's uncertainty 
was already proportionate to its error rate. A global $T < 1$ 
uniformly increases confidence across all classes; for SNII 
this produces overconfidence where none previously existed. 
A single global temperature parameter cannot simultaneously 
optimise calibration for classes with heterogeneous miscalibration 
characteristics. Class-specific temperature parameters would be 
required for a complete fix.

\begin{deluxetable}{lcccc}
\tablecaption{ALeRCE per-class ECE before and after temperature 
scaling ($T = 0.359$). \label{tab:tempscaling}}
\tablehead{
  \colhead{Class} & \colhead{$N$} & 
  \colhead{ECE (before)} & \colhead{ECE (after)} & 
  \colhead{$\Delta$ECE}
}
\startdata
SNIa  & 757 & 0.377 & 0.114 & $-0.263$ \\
SNII  & 560 & 0.099 & 0.189 & $+0.090$ \\
SNIbc & 204 & 0.360 & 0.108 & $-0.252$ \\
SLSN  &  55 & 0.225 & 0.140 & $-0.085$ \\
\enddata
\tablecomments{Positive $\Delta$ECE indicates worsening after 
temperature scaling. SNII is the only class for which global 
temperature scaling is counterproductive.}
\end{deluxetable}

\subsection{Fink Binary Classifiers}
\label{sec:results:fink}

Figure~\ref{fig:fink_reliability} shows reliability diagrams for 
Fink's two binary classifiers evaluated on 1,368 objects. Both 
exhibit severe miscalibration, but with qualitatively different 
failure modes.

The Random Forest classifier (\texttt{rf\_snia\_vs\_nonia}) assigns 
near-zero SNIa probability to virtually all objects: mean score 
$0.011$ against a true SNIa fraction of $0.471$. The reliability 
curve is essentially flat at the base rate regardless of predicted 
confidence, indicating the classifier has no discriminative power 
for SNIa identification in this sample 
($\mathrm{ECE} = 0.464$, 95\% CI: $[0.437, 0.490]$).

\begin{figure}
\centering
\includegraphics[width=\columnwidth]{../figures/fig3_fink_reliability}
\caption{Reliability diagrams for Fink's two binary classifiers 
evaluated on 1,368 ZTF transients. Left: Random Forest 
(\texttt{rf\_snia\_vs\_nonia}), ECE $= 0.464$. The flat curve 
indicates no discriminative power for SNIa identification. Right: 
SuperNNova (\texttt{snn\_snia\_vs\_nonia}), ECE $= 0.304$. Weak 
discriminative signal persists but is severely suppressed.}
\label{fig:fink_reliability}
\end{figure}

SuperNNova (\texttt{snn\_snia\_vs\_nonia}) shows qualitatively 
better behaviour: mean score $0.222$ against true fraction $0.471$, 
with a reliability curve that rises weakly with confidence 
($\mathrm{ECE} = 0.304$, 95\% CI: $[0.278, 0.332]$). The 
classifier retains some discriminative signal but suppresses it 
substantially.

Temperature scaling applied to both Fink classifiers hits the 
upper optimisation bound ($T \approx 10$) in both cases, 
achieving only modest ECE reductions: 24\% for RF 
($0.464 \to 0.353$) and 25\% for SuperNNova 
($0.304 \to 0.229$). When the optimal $T$ hits the bound, it 
indicates the score distribution lacks the monotonic structure 
temperature scaling requires. For Fink's RF classifier, no 
post-hoc transformation of the output scores can recover 
well-calibrated probabilities; the discriminative failure is 
structural.

\subsection{Cross-Broker Comparison}
\label{sec:results:comparison}

Table~\ref{tab:comparison} summarises calibration results across 
all three classifiers. ALeRCE's light curve classifier, despite 
exhibiting the systematic underconfidence described above, is 
substantially better calibrated than either Fink classifier. 
Critically, ALeRCE's miscalibration is correctable: temperature 
scaling reduces ECE to $0.031$, below the 0.05 significance 
threshold. Fink's classifiers, particularly the RF classifier, 
exhibit structural miscalibration that post-hoc methods cannot 
adequately address.

\begin{deluxetable}{lcccc}
\tablecaption{Calibration comparison across brokers and classifiers.
\label{tab:comparison}}
\tablehead{
  \colhead{Classifier} & \colhead{ECE} & 
  \colhead{95\% CI} & \colhead{ECE (post-TS)} &
  \colhead{$T$}
}
\startdata
ALeRCE LC    & 0.259 & [0.237, 0.280] & 0.031 & 0.359 \\
Fink RF      & 0.464 & [0.437, 0.490] & 0.353 & $\sim$10 \\
Fink SNN     & 0.304 & [0.278, 0.332] & 0.229 & $\sim$10 \\
\enddata
\tablecomments{Post-TS: ECE after temperature scaling. 
$T \sim 10$ indicates the optimiser hit the upper bound, 
signifying structural miscalibration not addressable by 
temperature scaling.}
\end{deluxetable}

These findings establish that miscalibration is not uniform 
across production astronomical brokers. The differences are 
scientifically consequential: a broker with ECE $= 0.464$ 
cannot support reliable deferral decisions, while a broker 
with ECE $= 0.031$ post-calibration approaches the reliability 
required for automated follow-up prioritisation.

\section{Discussion}
\label{sec:discussion}
\subsection{Systematic Underconfidence as a Structural Property}
\label{sec:discussion:underconfidence}

The dominant finding of this study is that ALeRCE's light curve 
classifier exhibits systematic underconfidence rather than the 
overconfidence commonly reported for deep learning classifiers 
\citep{Guo2017, Minderer2021}. This distinction matters for 
operational astronomy. Overconfident classifiers produce dangerous 
silent failures — high-confidence wrong predictions that are acted 
upon without human review. Underconfident classifiers produce 
unnecessary deferrals — objects the classifier effectively 
discriminates but assigns low confidence, causing them to be 
flagged for spectroscopic follow-up that is not scientifically 
required.

In the resource-constrained regime of LSST-era astronomy, where 
spectroscopic follow-up capacity is approximately $10^5$ 
observations per year against $10^7$ nightly alerts, unnecessary 
deferrals carry a direct opportunity cost. Every spectrum allocated 
to an object the classifier could have confidently classified 
correctly is a spectrum not allocated to a genuinely ambiguous 
object or a rare high-priority target. Quantifying and correcting 
underconfidence is therefore not merely a technical calibration 
exercise but a resource allocation problem.

\subsection{Why Random Forest Suppresses Confidence}
\label{sec:discussion:rf}

The underconfidence we observe is mechanistically expected for 
Random Forest classifiers in class-imbalanced settings. RF 
probability estimates are computed as the fraction of trees 
voting for each class. In a forest trained on an imbalanced 
dataset — SNIa comprises 47\% of our sample — the decision 
boundaries are dominated by the majority class, and minority 
class probabilities are systematically suppressed \citep{Niculescu-Mizil2005}. 
This suppression is strongest for classes the classifier 
discriminates most reliably, because reliable discrimination 
concentrates vote fractions near the boundary between correct 
and incorrect, rather than near 1.0.

This explains the inverse relationship we observe between accuracy 
and ECE across classes. SNIa and SNIbc are classified with 
accuracies above 0.83, yet their ECEs exceed 0.36. SNII, 
classified at only 0.495 accuracy, has ECE of 0.103. The 
classifier is most miscalibrated precisely where it has 
the most genuine discriminative signal.

The practical implication is that RF-based astronomical classifiers 
should not be used for confidence-based deferral decisions without 
post-hoc calibration. The raw probability outputs do not represent 
reliable confidence estimates, even when classification accuracy 
is high.

\subsection{Temperature Scaling: Efficacy and Limits}
\label{sec:discussion:tempscaling}

Temperature scaling with $T = 0.359$ reduces ALeRCE's ECE from 
$0.259$ to $0.031$ — an 88\% improvement with a single scalar 
parameter requiring no model retraining. This result is 
encouraging for operational deployment: existing broker 
infrastructure could be augmented with a post-hoc calibration 
layer at negligible computational cost.

However, the per-class analysis reveals an important limitation. 
Global temperature scaling worsens SNII calibration from $0.099$ 
to $0.189$, because SNII was already well-calibrated before 
scaling. A single $T$ cannot simultaneously optimise calibration 
for classes with heterogeneous miscalibration structure. 

Two approaches could address this. First, class-specific 
temperature parameters — fitting a separate $T_k$ for each 
class $k$ — would allow independent correction of each 
class's calibration curve. Second, more expressive post-hoc 
methods such as isotonic regression \citep{Zadrozny2002} or 
matrix scaling \citep{Kull2019} could model class interactions 
that scalar temperature cannot capture. We leave systematic 
comparison of these methods to future work, noting that the 
simple temperature scaling result already achieves practically 
useful calibration for the majority of classes.

\subsection{Fink's Structural Miscalibration}
\label{sec:discussion:fink}

The failure of temperature scaling for Fink's RF classifier — 
evidenced by the optimiser hitting the upper bound at $T \approx 10$ 
— indicates a qualitatively different problem from ALeRCE's 
underconfidence. When no finite temperature produces a well-calibrated 
output, the score distribution lacks the monotonic relationship 
between predicted probability and empirical accuracy that 
temperature scaling requires.

Inspection of the score distribution confirms this: 
\texttt{rf\_snia\_vs\_nonia} assigns scores below $0.05$ to 
over 90\% of objects, including the majority of true SNIa. 
The classifier is not underconfident in the same sense as 
ALeRCE — it is effectively non-discriminative for SNIa 
identification on this dataset. This may reflect a domain 
shift between Fink's training data and the BTS sample, or 
a fundamental limitation of the binary RF formulation for 
this classification problem.

SuperNNova retains weak discriminative structure — its 
reliability curve rises slightly with confidence — but 
insufficient for effective calibration. The deep learning 
architecture of SuperNNova provides some advantage over 
the RF approach, consistent with broader trends in the 
literature \citep{Moller2019, Villar2019}.

\subsection{Implications for Deferral-Aware Classification}
\label{sec:discussion:deferral}

This calibration study motivates the deferral-aware 
classification framework that constitutes the broader 
research programme of which this paper is the first component. 
Learning-to-defer frameworks \citep{Madras2018, Mozannar2020} 
require reliable confidence estimates as input to the deferral 
policy: an object should be deferred when classifier uncertainty 
genuinely reflects the difficulty of the classification, not 
when a well-discriminating classifier mechanically suppresses 
its own signal.

Our results suggest that ALeRCE, post-calibration, provides 
confidence estimates of sufficient quality to support deferral 
decisions ($\mathrm{ECE} = 0.031$). Fink's classifiers, in 
their current form, do not. This finding directly informs 
broker selection for downstream deferral framework development.

The TDE taxonomy gap identified in Section~\ref{sec:results:tde} 
represents a separate but related challenge. Deferral frameworks 
require that all scientifically relevant classes appear in the 
classifier's output space. A broker that cannot classify TDE 
by design will never defer TDE candidates for spectroscopic 
follow-up — it will instead misclassify them as the nearest 
available class with high confidence, producing exactly the 
silent failure mode that deferral is designed to prevent.

\subsection{Limitations}
\label{sec:discussion:limitations}

Several limitations of this study should be acknowledged. 
First, we evaluate calibration on a single broker snapshot; 
ALeRCE and Fink update their classifiers periodically, and 
our results reflect classifier versions active during our 
data collection period. Second, the BTS magnitude limit 
restricts our sample to bright, low-redshift transients; 
calibration properties may differ for fainter populations 
that will dominate LSST alert streams. Third, our temperature 
scaling analysis is conducted on the evaluation set itself, 
providing an optimistic estimate of post-hoc calibration 
performance; deployment would require fitting $T$ on a 
held-out validation set.

Finally, we note that ECE as defined in 
Equation~\ref{eq:ece} is sensitive to bin width and the 
distribution of samples across bins \citep{Nixon2019}. 
Our bootstrap confidence intervals partially address this 
by quantifying sampling uncertainty, but do not capture 
binning sensitivity. Adaptive binning methods \citep{Roelofs2022} 
may provide more robust ECE estimates for highly imbalanced 
confidence distributions such as those observed for Fink's 
RF classifier.

\section{Conclusions}
\label{sec:conclusions}
We have presented the first systematic calibration audit of production 
astronomical transient classifiers, evaluating ALeRCE's light curve 
classifier and two Fink classifiers against 1,606 spectroscopically 
confirmed ZTF transients from the Bright Transient Survey. Our results 
demonstrate that miscalibration is systematic, broker-dependent, and 
consequential for confidence-based follow-up prioritisation.

ALeRCE's classifier exhibits systematic underconfidence driven by 
Random Forest probability suppression under class imbalance 
($\mathrm{ECE} = 0.259$). This miscalibration is correctable: 
temperature scaling with $T = 0.359$ reduces ECE by 88\% to $0.031$, 
requiring no model retraining. The correction is not uniform across 
classes — SNII, the only well-calibrated class before scaling, 
worsens under a global temperature parameter — motivating class-aware 
calibration schemes for future work. TDE classification fails entirely 
due to a taxonomy gap in ALeRCE's output space, an operational problem 
requiring resolution before LSST first light. Fink's classifiers exhibit 
structural miscalibration that post-hoc methods cannot adequately address, 
with temperature scaling optimisation hitting the upper bound for both 
classifiers.

These findings establish three priorities for LSST-era broker 
development. First, calibration metrics — ECE, reliability diagrams, 
and per-class confidence intervals — should be reported alongside 
accuracy for all production classifiers as standard practice. Second, 
post-hoc calibration layers should be integrated into broker pipelines; 
the negligible computational cost of temperature scaling is justified 
by the substantial improvement in probability reliability. Third, 
taxonomy completeness must be verified against LSST science priorities 
before deployment. The calibrated ALeRCE outputs demonstrated here 
provide a suitable empirical foundation for the deferral-aware 
classification framework we develop in subsequent work, where reliable 
confidence estimates are prerequisite to learning when to involve 
human astronomers in follow-up decisions.

\bibliography{references}
\bibliographystyle{aasjournal}

\end{document}